% Holstein–Primakoff Expansion
% Converted on 2026-10-08 from docs/lswt/01-derivation/holstein-primakoff-expansion.md (status: draft, source-section: Spin to Boson Transformations).
% Review status: draft; awaiting the user's physics and mathematics review.

\hypertarget{sec-lswt-holstein-primakoff-expansion}{%
\section{Holstein--Primakoff Expansion}\label{sec-lswt-holstein-primakoff-expansion}}

The Holstein--Primakoff representation expresses fluctuations about a local spin direction in terms of boson operators. The local axes are defined in \lswtnoteref{LSWT foundations}{Classical Order and Local Frame}.

\subsection{Leading Spin Operators}

At the order used in linear spin-wave theory, the transverse spin operators are

\[
\hat{\widetilde S}_I^+
=\sqrt{2S_I}\,\hat a_I+O(S_I^{-1/2}),
\qquad
\hat{\widetilde S}_I^-
=\sqrt{2S_I}\,\hat a_I^\dagger+O(S_I^{-1/2}),
\]

and the longitudinal component is

\[
\hat{\widetilde S}_I^0=S_I-\hat n_I,
\qquad
\hat n_I=\hat a_I^\dagger\hat a_I.
\]

Here, \(S_I\) is the spin length, and \(\hat n_I\) counts the local Holstein--Primakoff bosons. The superscripts \(\pm\) denote the local raising and lowering components, while \(0\) denotes the longitudinal component. The boson occupation measures the reduction of that longitudinal spin component from \(S_I\).

\subsection{Fluctuations and Classical Stationarity}

The spin fluctuation separates into transverse and longitudinal parts relative to the local direction \(\mathbf e_I^0\):

\[
\delta\hat{\mathbf S}_I
=\delta\hat{\mathbf S}_I^\perp
 +\delta\hat{\mathbf S}_I^\parallel,
\qquad
\delta\hat{\mathbf S}_I^\parallel=-\hat n_I\mathbf e_I^0.
\]

When the reference configuration is an extremum of the classical energy under fixed-spin-length constraints, the linear boson terms vanish. With a Lagrange multiplier \(\lambda_I\) for the constraint at each site,

\[
\frac{\partial E_{\mathrm{cl}}}{\partial\mathbf S_I}
=\lambda_I\mathbf S_I,
\qquad
\frac{\partial E_{\mathrm{cl}}}{\partial\mathbf S_I}
\cdot\delta\hat{\mathbf S}_I^\perp=0.
\]

The energy gradient is parallel to the classical spin, whereas the leading transverse fluctuation is perpendicular to it. Their vanishing contraction removes the linear contribution. The quadratic terms generated by the expansion are collected in Real-Space Boson Hamiltonian (Section~\ref{sec-lswt-real-space-boson-hamiltonian}).

\subsection{References}

\subsubsection{Internal Documents}

\begin{itemize}
\tightlist
\item
  \lswtnoteref{LSWT foundations}{Classical Order and Local Frame}: defines the local spin frame and quantization axis used in the expansion.
\item
  Real-Space Boson Hamiltonian (Section~\ref{sec-lswt-real-space-boson-hamiltonian}): collects the quadratic terms generated by the expansion.
\item
  \lswtnoteref{LSWT foundations}{Notation and Conventions}: defines the HP-boson symbols and site indices.
\end{itemize}

\subsubsection{External Sources}

\begin{itemize}
\tightlist
\item
  T. Holstein and H. Primakoff, ``Field Dependence of the Intrinsic Domain Magnetization of a Ferromagnet,'' \emph{Physical Review} \textbf{58}, 1098-1113 (1940), \href{https://doi.org/10.1103/PhysRev.58.1098}{doi:10.1103/PhysRev.58.1098}: original source for the Holstein--Primakoff boson representation.
\end{itemize}
